%% GENERATED by tools/gen_latex.py -- edit content/ch03.py instead.
%% Build:  latexmk -pdf ch03.tex     (or: pdflatex ch03.tex, twice)
\documentclass[aspectratio=169,11pt,t]{beamer}
\usepackage{itbpro}

\coursetitle{Introduction to TensorFlow, PyTorch, JAX, and Keras}
\coursesubtitle{The same Dense layer written four times --- so the design differences between the frameworks are something you can see rather than something you are told.}
\coursekicker{Chapter 3}
\coursesource{Chollet \& Watson, Deep Learning with Python 3e --- chapter 3}
\courseduration{2.5 hours}
\coursepresenter{Rahman Indra Kesuma, S.Kom., M.Cs.}{Teaching Assistant}
\coursebrand{AI for Professional \textperiodcentered\ ITB \textperiodcentered\ Ch 3}

\title{Introduction to TensorFlow, PyTorch, JAX, and Keras}
\author{Rahman Indra Kesuma, S.Kom., M.Cs.}
\date{}

\begin{document}
\covertitle

\framekicker{Learning outcomes}
\begin{frame}[shrink=25]{Session Objectives}
\leadpar{By the end of this session, participants can:}
\begin{isteps}
\item Name the \textbf{three capabilities} every modern framework provides, and what distinguishes them beyond that.
\item Compute gradients with \textbf{GradientTape, \inlinecode{.backward()}, and \inlinecode{jax.grad()}}, and explain why the three look so different.
\item Distinguish \textbf{stateful imperative} (TF, PyTorch) from \textbf{stateless functional} (JAX), and say what that costs when writing a training loop.
\item Switch the Keras 3 backend without changing a line of model code.
\item Write a custom \inlinecode{Layer} with \inlinecode{build()} and \inlinecode{call()}, then drive it through \inlinecode{compile()} and \inlinecode{fit()}.
\end{isteps}

\vskip 6pt
\vskip 4pt
\reslink{Book}{Chapter 3 --- full text}{https://deeplearningwithpython.io/chapters/chapter03\_introduction-to-ml-frameworks}
\reslink{Notebook}{01\_three\_frameworks\_side\_by\_side.ipynb}{../../notebooks/ch03/01\_three\_frameworks\_side\_by\_side.ipynb}
\reslink{Notebook}{02\_keras3\_switch\_backend.ipynb}{../../notebooks/ch03/02\_keras3\_switch\_backend.ipynb}
\reslink{Notebook}{03\_custom\_layer\_and\_fit.ipynb}{../../notebooks/ch03/03\_custom\_layer\_and\_fit.ipynb}
\reslink{Repo}{Author's official notebook}{https://github.com/fchollet/deep-learning-with-python-notebooks/blob/master/chapter03\_introduction-to-ml-frameworks.ipynb}
\vskip 2pt
\end{frame}

\framekicker{Section 3.1}
\begin{frame}[shrink=25]{How the field arrived here}
\begin{center}
\includegraphics[width=0.92\linewidth,height=0.60\textheight,keepaspectratio]{figs/ch03-timeline.tex.pdf}
\end{center}
\vskip 3pt{\color{ink3}\fontsize{7.4}{9.4}\selectfont Automatic differentiation is from 1964. What was new in 2009 was combining it with GPU computation.\par}
\begin{iband}[signal]
One number worth keeping: by \textbf{mid-2016, more than half} of TensorFlow's users reached it \hilite{through Keras}
\end{iband}
\end{frame}

\framekicker{Section 3.1}
\begin{frame}[shrink=25]{Three capabilities every one of them has}
\begin{minipage}[t]{0.3253\textwidth}
\begin{icardaccent}
\cardh{Automatic differentiation}
For any differentiable function you write, not just for a fixed catalogue of layers.
\end{icardaccent}
\end{minipage}%
\hfill
\begin{minipage}[t]{0.3253\textwidth}
\begin{icardaccent}
\cardh{Tensor computation}
On CPUs, GPUs, and specialised hardware such as TPUs.
\end{icardaccent}
\end{minipage}%
\hfill
\begin{minipage}[t]{0.3253\textwidth}
\begin{icardaccent}
\cardh{Distributed execution}
Across devices and across machines.
\end{icardaccent}
\end{minipage}%
\begin{iband}[signal]
Because all of them provide all three, choosing a framework is \hilite{not a question of capability}. It is a question of writing style, ecosystem, and speed.
\end{iband}
\end{frame}

\framekicker{Section 3.2}
\begin{frame}[shrink=25]{Keras on top, three engines underneath}
\begin{center}
\includegraphics[width=0.92\linewidth,height=0.60\textheight,keepaspectratio]{figs/ch03-stack.tex.pdf}
\end{center}
\vskip 3pt{\color{ink3}\fontsize{7.4}{9.4}\selectfont Keras needs a backend. NumPy can be plugged in but cannot train --- it has no gradient API.\par}
\end{frame}

\framekicker{Section 3.2}
\begin{frame}[shrink=25]{Where the line between them falls}
\begin{iquote}
\quotetext{Keras is like a prefabricated building kit, while TensorFlow, PyTorch, and JAX are like raw materials used in construction.}\quotecite{Chollet \& Watson, section 3.2}
\end{iquote}
\begin{minipage}[t]{0.485\textwidth}
{\fontsize{9.4}{12.6}\selectfont\color{fg2}\textbf{Low level} --- tensors, tensor operations, and backpropagation.\par}\vskip 4pt

\end{minipage}%
\hfill
\begin{minipage}[t]{0.485\textwidth}
{\fontsize{9.4}{12.6}\selectfont\color{fg2}\textbf{High level} --- layers combined into models, losses, optimizers, metrics, and training loops.\par}\vskip 4pt

\end{minipage}%
\end{frame}

\coursesection{01}{TensorFlow}{Immutable tensors, Variables for state, a tape for gradients.}

\framekicker{Section 3.3}
\begin{frame}[shrink=25]{Constants are immutable; state needs a Variable}
{\fontsize{9.4}{12.6}\selectfont\color{fg2}This is the first thing that surprises people coming from NumPy: a TensorFlow tensor cannot be assigned into. Anything that has to change during training must be a \inlinecode{tf.Variable}.\par}\vskip 4pt
\icode{python}{tensors and variables}{listings/ch03-01.py}
\begin{iband}[amber]
Forgetting this produces a confusing error at the worst moment: \hilite{your weights simply refuse to update}
\end{iband}
\end{frame}

\framekicker{Section 3.3}
\begin{frame}[shrink=25]{Operations, and one naming detail to notice}
{\fontsize{9.4}{12.6}\selectfont\color{fg2}The operation names mostly follow NumPy, with occasional divergences. Watch the keyword in the last line.\par}\vskip 4pt
\icode{python}{tensor operations}{listings/ch03-02.py}
{\fontsize{9.4}{12.6}\selectfont\color{fg2}That \inlinecode{axis} keyword becomes \inlinecode{dim} in PyTorch --- a small difference that costs real time when porting code.\par}\vskip 4pt
\end{frame}

\framekicker{Section 3.3}
\begin{frame}[shrink=25]{GradientTape: record, then ask}
{\fontsize{9.4}{12.6}\selectfont\color{fg2}TensorFlow records the operations performed inside a tape's scope, then replays them backwards when you ask for a gradient.\par}\vskip 4pt
\icode{python}{gradients}{listings/ch03-03.py}
\begin{iband}[signal]
Variables are watched automatically because they are what you normally differentiate. \hilite{Constants must be opted in}
\end{iband}
\end{frame}

\framekicker{Section 3.3}
\begin{frame}[shrink=25]{Compilation: graph mode and XLA}
{\fontsize{9.4}{12.6}\selectfont\color{fg2}Decorating a function replaces it with a compiled program. The first call is slower; every call after that is faster.\par}\vskip 4pt
\icode{python}{two levels of compilation}{listings/ch03-04.py}
{\fontsize{9.4}{12.6}\selectfont\color{fg2}This is the same mechanism that later lets a model be \textbf{exported without Python} --- chapter 6 uses it for deployment.\par}\vskip 4pt
\end{frame}

\framekicker{Section 3.3}
\begin{frame}[shrink=25]{TensorFlow: where it wins and where it hurts}
\begin{minipage}[t]{0.494\textwidth}
\begin{icardgood}
\cardh{Strengths}
Fast via graph mode and XLA $\cdot$ very feature-complete (string tensors, ragged tensors) $\cdot$ outstanding \inlinecode{tf.data} for preprocessing $\cdot$ \textbf{the most mature ecosystem for production, mobile, and browser deployment}.
\end{icardgood}
\end{minipage}%
\hfill
\begin{minipage}[t]{0.494\textwidth}
\begin{icardbad}
\cardh{Weaknesses}
A sprawling API with thousands of operations $\cdot$ inconsistent with NumPy in places $\cdot$ less support on Hugging Face than PyTorch.
\end{icardbad}
\end{minipage}%
{\fontsize{9.4}{12.6}\selectfont\color{fg2}For teams that must deploy on-premise, that production ecosystem is usually the argument that decides it.\par}\vskip 4pt
\end{frame}

\framekicker{Section 3.3}
\begin{frame}[shrink=25]{A whole training step in TensorFlow}
{\fontsize{9.4}{12.6}\selectfont\color{fg2}Putting the pieces together: a compiled step function that runs the forward pass under a tape, reads the gradients, and updates two variables in place.\par}\vskip 4pt
\icode{python}{an end-to-end training step}{listings/ch03-05.py}
\begin{iband}[signal]
\inlinecode{assign\_sub} mutates the variable. Hold that image --- \hilite{the JAX version of this same loop cannot mutate anything}, and the difference is instructive.
\end{iband}
\end{frame}

\coursesection{02}{PyTorch}{Mutable tensors, .backward() fills .grad, eager by default.}

\framekicker{Section 3.4}
\begin{frame}[shrink=25]{Tensors you can assign into}
{\fontsize{9.4}{12.6}\selectfont\color{fg2}PyTorch takes the opposite decision from TensorFlow: its tensors behave like NumPy arrays and can be written to directly.\par}\vskip 4pt
\icode{python}{tensors and parameters}{listings/ch03-06.py}
{\fontsize{9.4}{12.6}\selectfont\color{fg2}\inlinecode{Parameter} does not change the maths; it marks a tensor as something an optimizer should own.\par}\vskip 4pt
\end{frame}

\framekicker{Section 3.4}
\begin{frame}[shrink=25]{No tape --- the graph is built as you go}
{\fontsize{9.4}{12.6}\selectfont\color{fg2}Each forward pass builds a one-time computation graph. Calling \inlinecode{.backward()} on a scalar walks it in reverse and fills in a \inlinecode{.grad} on every tensor involved.\par}\vskip 4pt
\icode{python}{gradients, and the trap}{listings/ch03-07.py}
\begin{iband}[rose]
That last line is the classic PyTorch bug. The next \inlinecode{.backward()} \textbf{adds} to the existing gradient rather than replacing it, so \hilite{training silently stops working}
\end{iband}
\end{frame}

\framekicker{Section 3.4}
\begin{frame}[shrink=25]{The three-line incantation}
{\fontsize{9.4}{12.6}\selectfont\color{fg2}Model, loss, and optimizer are aware of each other, and a training step is always the same three calls in the same order.\par}\vskip 4pt
\icode{python}{the PyTorch training pattern}{listings/ch03-08.py}
{\fontsize{9.4}{12.6}\selectfont\color{fg2}Learn those three lines in that order and most PyTorch code becomes readable.\par}\vskip 4pt
\end{frame}

\framekicker{Section 3.4}
\begin{frame}[shrink=25]{PyTorch: where it wins and where it hurts}
\begin{minipage}[t]{0.494\textwidth}
\begin{icardgood}
\cardh{Strengths}
Eager by default, which makes debugging the easiest of the three $\cdot$ \textbf{first-class support on Hugging Face}, and that is the single biggest driver of its adoption.
\end{icardgood}
\end{minipage}%
\hfill
\begin{minipage}[t]{0.494\textwidth}
\begin{icardbad}
\cardh{Weaknesses}
Internally inconsistent API (\inlinecode{axis} sometimes becomes \inlinecode{dim}) $\cdot$ in the authors' assessment \textbf{the slowest of the major frameworks} $\cdot$ \inlinecode{torch.compile()} is still full of edge cases and little used.
\end{icardbad}
\end{minipage}%
{\fontsize{9.4}{12.6}\selectfont\color{fg2}If your team already lives on Hugging Face, PyTorch will feel like home. That is a legitimate reason to choose it.\par}\vskip 4pt
\end{frame}

\framekicker{Section 3.4}
\begin{frame}[shrink=25]{Compilation in PyTorch, and why it is rarely used}
{\fontsize{9.4}{12.6}\selectfont\color{fg2}PyTorch gained a compiler comparatively late. It can be applied to a model or used as a decorator.\par}\vskip 4pt
\icode{python}{torch.compile}{listings/ch03-09.py}
\begin{iband}[amber]
Unlike TensorFlow and JAX, \textbf{most PyTorch code runs eagerly, uncompiled}. The book is blunt that the Dynamo compiler is still full of edge cases and that \hilite{only a small percentage of users employ it}
\end{iband}
\end{frame}

\coursesection{03}{JAX}{Stateless functions. Gradients as a transformation of a function.}

\framekicker{Section 3.5}
\begin{frame}[shrink=25]{Stateless --- including the random numbers}
{\fontsize{9.4}{12.6}\selectfont\color{fg2}JAX functions keep no state between calls. State is passed in as arguments and returned as results --- and that applies to the random number generator too.\par}\vskip 4pt
\icode{python}{arrays, keys, updates}{listings/ch03-10.py}
\begin{iband}[signal]
It reads as extra work, and it is. What you buy is computation that \hilite{parallelises automatically without synchronisation}, and results that are exactly reproducible.
\end{iband}
\end{frame}

\framekicker{Section 3.5}
\begin{frame}[shrink=25]{jax.grad(): a function in, a function out}
{\fontsize{9.4}{12.6}\selectfont\color{fg2}JAX does not record operations. It \textbf{transforms your function} into a different function that computes gradients.\par}\vskip 4pt
\icode{python}{the transformation}{listings/ch03-11.py}
\begin{itable}{>{\raggedright\arraybackslash}p{0.3948\textwidth}>{\raggedright\arraybackslash}p{0.5452\textwidth}}
\thead{What you need} & \thead{What you call} \\ \midrule
Gradients only & \inlinecode{jax.grad(f)} \\
Loss \textbf{and} gradients together & \inlinecode{jax.value\_and\_grad(f)} --- cheaper \\
Plus auxiliary outputs & \inlinecode{jax.value\_and\_grad(f, has\_aux=True)} \\
\end{itable}
\end{frame}

\framekicker{Section 3.5}
\begin{frame}[shrink=25]{What a full training step looks like}
{\fontsize{9.4}{12.6}\selectfont\color{fg2}Because nothing mutates, the step function has to \textbf{return} the new weights. That single constraint shapes all JAX code.\par}\vskip 4pt
\icode{python}{a JAX training step}{listings/ch03-12.py}
\begin{iband}[amber]
Notice the loop: \inlinecode{W} and \inlinecode{b} are \hilite{threaded through by hand}. Nothing is updated in place, anywhere.
\end{iband}
\end{frame}

\framekicker{Section 3.5}
\begin{frame}[shrink=25]{JAX: where it wins and where it hurts}
\begin{minipage}[t]{0.494\textwidth}
\begin{icardgood}
\cardh{Strengths}
In the authors' assessment \textbf{the fastest of the three} $\cdot$ perfect NumPy API consistency, so no surprises $\cdot$ built for XLA and TPUs from the beginning.
\end{icardgood}
\end{minipage}%
\hfill
\begin{minipage}[t]{0.494\textwidth}
\begin{icardbad}
\cardh{Weaknesses}
Metaprogramming plus compilation makes debugging \textbf{markedly harder} than eager execution $\cdot$ low-level training loops are more verbose than TF or PyTorch.
\end{icardbad}
\end{minipage}%
\end{frame}

\framekicker{Sections 3.3 -- 3.5}
\begin{frame}[shrink=25]{The same job, three ways of asking for it}
\begin{center}

\includegraphics[width=0.92\linewidth,height=0.60\textheight,keepaspectratio]{figs/ch03-gradstyle.tex.pdf}
\end{center}
\vskip 3pt{\color{ink3}\fontsize{7.4}{9.4}\selectfont The gradient is the same quantity in all three; what differs is who holds the state while it is computed.\par}
\end{frame}

\framekicker{Sections 3.3 -- 3.5}
\begin{frame}[shrink=25]{Side by side}
\begin{itable}{>{\raggedright\arraybackslash}p{0.1598\textwidth}>{\raggedright\arraybackslash}p{0.2632\textwidth}>{\raggedright\arraybackslash}p{0.2632\textwidth}>{\raggedright\arraybackslash}p{0.2538\textwidth}}
\thead{} & \thead{TensorFlow} & \thead{PyTorch} & \thead{JAX} \\ \midrule
\textbf{Paradigm} & Stateful imperative & Stateful imperative & Stateless functional \\
\textbf{Tensors} & Immutable (\inlinecode{Variable} for state) & Mutable & Immutable (\inlinecode{.at[].set()}) \\
\textbf{Gradients} & \inlinecode{GradientTape} & \inlinecode{.backward()} $\rightarrow$ \inlinecode{.grad} & \inlinecode{jax.grad()} transformation \\
\textbf{Compilation} & \inlinecode{@tf.function}, XLA & \inlinecode{@torch.compile} (Dynamo) & \inlinecode{@jax.jit} (XLA) \\
\textbf{Debugging} & Harder in graph mode & Easiest & Hardest (functional + JIT) \\
\textbf{Ecosystem} & Production tooling & Hugging Face, research & Research, Google scale \\
\end{itable}
\end{frame}

\framekicker{Sections 3.3 -- 3.5}
\begin{frame}[shrink=25]{About those speed claims}
\begin{iband}[amber]
The book's ranking --- JAX fastest, PyTorch slowest, a 20--30\% spread and up to 3--5$\times$ on large models --- is the \textbf{authors' assessment}, \hilite{not a benchmark this course ran}
\end{iband}
{\fontsize{9.4}{12.6}\selectfont\color{fg2}Treat it as a direction to investigate, and measure it yourself on your own workload before it becomes a procurement argument.\par}\vskip 4pt
\end{frame}
\note{Do not let the room turn this into a framework war. The point is that real design differences exist and Keras papers over them.}

\coursesection{04}{Keras}{One model definition, three engines behind it.}

\framekicker{Section 3.6}
\begin{frame}[shrink=25]{Switching backend without touching the model}
{\fontsize{9.4}{12.6}\selectfont\color{fg2}The backend is a configuration choice, not a code change. There are two ways to make it.\par}\vskip 4pt
\begin{minipage}[t]{0.485\textwidth}
\icode{python}{way 1 --- environment variable}{listings/ch03-13.py}
\ioutput{listings/ch03-14.txt}

\end{minipage}%
\hfill
\begin{minipage}[t]{0.485\textwidth}
\icode{json}{way 2 --- \textasciitilde{}/.keras/keras.json}{listings/ch03-15.json}

\end{minipage}%
\begin{iband}[rose]
The ordering is not negotiable: \inlinecode{os.environ[...]} must run \hilite{before the first \inlinecode{import keras}}. After Keras is imported, changing it has no effect --- and this is the number one confusion in the lab sessions.
\end{iband}
\end{frame}

\framekicker{Section 3.6.1}
\begin{frame}[shrink=25]{Layer: the unit everything is built from}
{\fontsize{9.4}{12.6}\selectfont\color{fg2}Every Keras component either is a \inlinecode{Layer} or works closely with one. A custom layer needs two methods: \inlinecode{build()} creates the weights, \inlinecode{call()} does the work.\par}\vskip 4pt
\icode{python}{listing 3.22 --- a custom Layer}{listings/ch03-16.py}
{\fontsize{9.4}{12.6}\selectfont\color{fg2}Note that \inlinecode{build()} receives the input shape. The layer does not need to be told it in advance.\par}\vskip 4pt
\end{frame}

\framekicker{Section 3.6.2}
\begin{frame}[shrink=25]{Automatic shape inference, and what it buys}
\begin{center}
\includegraphics[width=0.92\linewidth,height=0.60\textheight,keepaspectratio]{figs/ch03-layerlife.tex.pdf}
\end{center}
\vskip 3pt{\color{ink3}\fontsize{7.4}{9.4}\selectfont Weights are created on the first call, not at construction time.\par}
{\fontsize{9.4}{12.6}\selectfont\color{fg2}This is why \inlinecode{Sequential} in chapter 2 only had to name the number of units. There was \hilite{no \inlinecode{input\_shape} anywhere}, and now you know why.\par}\vskip 4pt
\end{frame}

\framekicker{Section 3.6.2}
\begin{frame}[shrink=25]{Seeing it happen}
{\fontsize{9.4}{12.6}\selectfont\color{fg2}Instantiate the layer, call it on data, and read back the output shape it inferred.\par}\vskip 4pt
\icode{python}{using the custom layer}{listings/ch03-17.py}
\ioutput{listings/ch03-18.txt}
{\fontsize{9.4}{12.6}\selectfont\color{fg2}The 784 was never written down. It was read off the data at the moment the layer was first used.\par}\vskip 4pt
\end{frame}

\framekicker{Section 3.6}
\begin{frame}[shrink=25]{Which backend, in practice}
\begin{center}
\includegraphics[width=0.92\linewidth,height=0.60\textheight,keepaspectratio]{figs/ch03-choose.tex.pdf}
\end{center}
\vskip 3pt{\color{ink3}\fontsize{7.4}{9.4}\selectfont A decision aid, not a law. The point of Keras is that this choice stays reversible.\par}
{\fontsize{9.4}{12.6}\selectfont\color{fg2}The authors suggest \textbf{JAX} for best performance --- but the same Keras code runs on all three, so the decision can be deferred and revisited.\par}\vskip 4pt
\end{frame}

\framekicker{Section 3.6.3}
\begin{frame}[shrink=25]{Models are graphs of layers --- and a choice of hypothesis space}
\begin{iquote}
\quotetext{The topology of a model defines a hypothesis space. By choosing a network topology, you constrain your space of possibilities to a specific series of tensor operations.}\quotecite{Chollet \& Watson, section 3.6.3}
\end{iquote}
{\fontsize{9.4}{12.6}\selectfont\color{fg2}That phrase --- \textbf{hypothesis space} --- is the one from chapter 1. Choosing an architecture is choosing what the model is allowed to learn. Chapter 7 covers the three ways of expressing that choice.\par}\vskip 4pt
\end{frame}

\framekicker{Section 3.6.4}
\begin{frame}[shrink=25]{compile(): three decisions, spelled out}
{\fontsize{9.4}{12.6}\selectfont\color{fg2}The string shorthand is convenient, but the object form is what you need as soon as a learning rate has to be tuned.\par}\vskip 4pt
\icode{python}{listing 3.26 --- two equivalent forms}{listings/ch03-19.py}
{\fontsize{9.4}{12.6}\selectfont\color{fg2}Built in and ready to use: \textbf{SGD, RMSprop, Adam}; losses including the crossentropies and MSE; metrics including accuracy, AUC, precision, and recall.\par}\vskip 4pt
\end{frame}

\framekicker{Section 3.6.5}
\begin{frame}[shrink=25]{fit(), and why validation data is not optional}
{\fontsize{9.4}{12.6}\selectfont\color{fg2}\inlinecode{fit()} runs the loop and returns a \inlinecode{History} object whose \inlinecode{.history} dictionary holds the per-epoch values of every metric --- which is what you plot.\par}\vskip 4pt
\icode{python}{listing 3.29 --- training with validation}{listings/ch03-20.py}
\begin{iquote}
\quotetext{The goal of machine learning is not to obtain models that perform well on the training data \ldots{} it is to obtain models that perform well in general, particularly on data points that the model has never encountered before.}\quotecite{Chollet \& Watson, section 3.6.5}
\end{iquote}
\end{frame}

\framekicker{Section 3.6.5}
\begin{frame}[shrink=25]{Holding data back, by hand}
\begin{center}
\includegraphics[width=0.92\linewidth,height=0.60\textheight,keepaspectratio]{figs/ch03-valsplit.tex.pdf}
\end{center}
\vskip 3pt{\color{ink3}\fontsize{7.4}{9.4}\selectfont Shuffle first, then split. Splitting an ordered array is one of the classic ways to get a meaningless validation score.\par}
\end{frame}

\framekicker{Section 3.6.5}
\begin{frame}[shrink=25]{\ldots{}and what that looks like in code}
{\fontsize{9.4}{12.6}\selectfont\color{fg2}\inlinecode{fit()} can do this for you with \inlinecode{validation\_split}, but doing it by hand once makes the shuffle step impossible to forget later.\par}\vskip 4pt
\icode{python}{listing 3.28 --- a manual split}{listings/ch03-21.py}
\begin{iband}[signal]
Chapter 5 shows what goes wrong when the shuffle is skipped: you can end up training on classes 0--7 and testing on 8--9, and \hilite{the code will not complain}
\end{iband}
\end{frame}

\framekicker{Section 3.6.6}
\begin{frame}[shrink=25]{Two ways to run inference, and when each is right}
{\fontsize{9.4}{12.6}\selectfont\color{fg2}Calling the model directly is fine for a small batch you already hold in memory. \inlinecode{predict()} is what you want for anything larger.\par}\vskip 4pt
\icode{python}{inference}{listings/ch03-22.py}
\begin{iband}[signal]
\inlinecode{predict()} iterates in batches, so it \hilite{will not exhaust memory} on a large array, and it hands back NumPy rather than backend tensors.
\end{iband}
\end{frame}

\framekicker{Section 3.7}
\begin{frame}[shrink=25]{Why this stack is a safe thing to learn}
\begin{ibullets}
\item \textbf{Python has won} the ML and data science ecosystem; the authors see no replacement within fifteen years.
\item All four frameworks are \textbf{stable}. New ones may appear, but are unlikely to displace existing workflows.
\item New hardware --- AMD GPUs and others --- must integrate with the existing frameworks, so it \hilite{does not disrupt your code}.
\end{ibullets}
\begin{iband}[signal]
Which is the actual argument for Keras: it has provided \textbf{future-proof stability since 2015}, and a pluggable backend is what lets it keep doing so.
\end{iband}
\end{frame}

\framekicker{Section 3.6.4}
\begin{frame}[shrink=25]{Metrics are watched; loss is optimised}
\begin{iband}[amber]
A \textbf{metric} is monitored but \hilite{never optimised for}. Only the \textbf{loss} is minimised. Confusing the two is an expensive mistake, and chapters 5 and 6 return to exactly why.
\end{iband}
{\fontsize{9.4}{12.6}\selectfont\color{fg2}The short version: many of the things you actually care about --- ROC AUC, for instance --- cannot be used as a loss because they are not differentiable.\par}\vskip 4pt
\end{frame}

\framekicker{Summary}
\begin{frame}[shrink=25]{What has to survive this chapter}
\begin{isteps}
\item Every major framework gives you \textbf{autodiff, GPU/TPU tensor computation, and distributed execution}. The rest is style.
\item \textbf{TensorFlow} --- immutable tensors plus \inlinecode{Variable}, \inlinecode{GradientTape}, and the strongest production ecosystem.
\item \textbf{PyTorch} --- mutable tensors, \inlinecode{.backward()}, eager, king of Hugging Face. Never forget \inlinecode{zero\_grad()}.
\item \textbf{JAX} --- stateless functional, \inlinecode{jax.grad()} as a transformation, explicit random keys, fastest by the authors' account.
\item \textbf{Keras 3} sits on all three. Set \inlinecode{KERAS\_BACKEND} \hilite{before} \inlinecode{import keras}.
\item A \inlinecode{Layer} with \inlinecode{build()} + \inlinecode{call()} is the unit everything is made of, and input shapes are inferred on first use.
\end{isteps}
\vskip 4pt
\reslink{NOTEBOOK}{01\_three\_frameworks\_side\_by\_side.ipynb}{../../course-slides/notebooks/ch03/01\_three\_frameworks\_side\_by\_side.ipynb}
\reslink{NEXT}{Chapter 4 --- Classification and regression}{../ch04/index.html}
\vskip 2pt
\end{frame}

\end{document}
